% Source-specific revision, 12 September 2026. Compile twice with XeLaTeX.
% Historical Eulerian near-cubic proposal preserved in proposal_support/archive/.
\documentclass[12pt,a4paper]{article}
\usepackage[margin=1in]{geometry}
\usepackage{fontspec}
\setmainfont{Times New Roman}
\usepackage{amsmath,amssymb,setspace,enumitem,titlesec}
\usepackage[hidelinks]{hyperref}
\doublespacing
\titleformat{\section}{\normalfont\normalsize\bfseries}{\thesection.}{0.5em}{}
\titlespacing*{\section}{0pt}{6pt}{3pt}
\titleformat{\subsection}{\normalfont\normalsize\bfseries}{\thesubsection.}{0.5em}{}
\titlespacing*{\subsection}{0pt}{6pt}{3pt}
\setlist[enumerate]{leftmargin=1.6em,itemsep=0pt,parsep=0pt,topsep=2pt}
\setlength{\parindent}{0pt}
\setlength{\parskip}{3pt}
\setlength{\emergencystretch}{2em}
\hypersetup{pdftitle={Audited Lie Reduction and Semi-Analytical PINNs for a Lagrangian Shallow-Water Master Equation},pdfauthor={Carl}}
\input{proposal_support/lagrange_evidence/pilot_values.tex}
\begin{document}

\begin{center}
\textbf{Audited Lie Reduction and Semi-Analytical PINNs\\
for a Lagrangian Shallow-Water Master Equation}\\
Carl \quad WKU MATH 4890\\
Proposed supervisor: Dr. Puneet Rana
\end{center}

\section{Research Problem}
This project takes Paliathanasis, Leon and Leach's 2026 paper, \emph{Lie Symmetries and Similarity Solutions for a Shallow-Water Model with Bed Elevation in Lagrange Variables}, as its primary source \cite{lagrange}. Its Equation (7) is a third-order master PDE for the particle-position map $\phi(t,s)$, where $t$ is time and $s$ labels fluid mass. Lie invariants compress $(t,s)$ into one similarity coordinate, leaving an ODE for a profile.

The paper supplies symmetry classifications and similarity families, including an extra scaling symmetry when the sediment-flux exponent is $m=3$. These are established starting points, not proposed discoveries. The research question is: \textbf{can independently checked reductions support reliable Taylor--PINN profile approximations, with explicit regularity, error and reconstruction checks?}

A source audit identifies a bed-transport term missing from the printed coordinate conversion under the usual pullback definition. The immediate target is therefore the \emph{printed master equation}; equivalence with the original Eulerian fluid system is a separate check, not an assumption. General boundary data, shocks and full-PDE well-posedness remain outside scope.

\section{Objectives}
\begin{enumerate}
\item Audit the model conversion and rederive the selected cubic reduction.
\item Establish regular-branch conditions and a local residual-to-profile error bound.
\item Compare Taylor, classical numerical and reduced PINN approximations using exact controls, perturbed profile data and reproducible cost/error measurements.
\end{enumerate}

\newpage
\section{Methodology}
\subsection{Source Model and Reconstruction Audit}
Using dimensionless variables with fixed $g,\beta>0$, the paper's Equation (7) is
\begin{equation}
\mathcal E_m[\phi]=
\partial_t\left(\phi_s\phi_{tt}-g\frac{\phi_{ss}}{\phi_s^2}\right)
-g\beta\,\partial_s\left(\phi_t^{m-1}\frac{\phi_{ts}}{\phi_s}\right)=0.
\label{eq:master}
\end{equation}
The source defines $\beta=mA_g\xi$, where $A_g$ is the sediment-flux coefficient and $\xi$ is the porosity factor. Set $m=3$ in the main study. Require $\phi_s>0$ for a locally invertible mass-to-position map, and $\phi_t>0$ for the selected positive-velocity branch.

For the physical bed field $B(t,x)$, define $b(t,s)=B(t,\phi(t,s))$. Then
\[
u=\phi_t,\quad h=\phi_s^{-1},\qquad
\partial_x=\phi_s^{-1}\partial_s,\quad
(\partial_t)_x=(\partial_t)_s-\frac{\phi_t}{\phi_s}\partial_s.
\]
Thus the Eulerian Exner equation transforms to
\begin{equation}
b_t-\frac{\phi_t}{\phi_s}b_s+
\beta\phi_t^{m-1}\frac{\phi_{ts}}{\phi_s}=0.
\label{eq:physicalbed}
\end{equation}
The transport term in \eqref{eq:physicalbed} is absent from the source's printed Equation (6). The master equation instead expresses compatibility of the \emph{printed} pair
\[
b_s=-A/g,\qquad b_t=-\beta\phi_t^{m-1}\phi_{ts}/\phi_s,\qquad
A=\phi_s\phi_{tt}-g\phi_{ss}/\phi_s^2.
\]
On a simply connected regular region, \eqref{eq:master} makes this pair integrable up to a bed-level constant. It does not imply \eqref{eq:physicalbed}.

The audit will retain both residuals. Correcting the transport model would require new symmetry and reduction checks; no classification or neural result will be transferred automatically to that different equation.

\newpage
\subsection{Cubic Similarity Reduction and Regularity}
The primary source's Proposition 1 and Section 3.5 supply
\[
X_5=t\partial_t-2s\partial_s,\qquad
\sigma=st^2,\qquad \phi(t,s)=F(\sigma).
\]
The flow sends $(t,s)$ to $(e^\epsilon t,e^{-2\epsilon}s)$, leaving $\sigma$ fixed. The invariant-surface condition $t\phi_t-2s\phi_s=0$ therefore reduces two independent variables to one \cite{olver}. Direct substitution into \eqref{eq:master}, independently of the printed expanded ODEs, gives
\begin{equation}
\mathcal E_3[F(st^2)]=t^{-1}R[F],\qquad
R[F]=a(\sigma,r)F'''+b(\sigma,r,v),
\label{eq:lift}
\end{equation}
where $r=F'$, $v=F''$, and
\[
\begin{aligned}
a&=8(1-g\beta)\sigma^3r-2g\sigma/r^2,\\
b&=4\sigma(1-4g\beta)r^2+8\sigma^2(3-5g\beta)rv\\
 &\quad+\left[8\sigma^3(1-g\beta)+4g\sigma/r^3\right]v^2.
\end{aligned}
\]
On a region with $\sigma>0$, $r\ge r_{\min}>0$ and $|a|\ge a_{\min}>0$, this is the regular IVP
\[
Y'=(r,v,-b/a)=G(\sigma,Y),\qquad Y=(F,r,v).
\]
The conditions prevent loss of map invertibility and division by a vanishing highest-derivative coefficient. Rank drops are recorded as unsupported continuation, not removed by cancellation.

\textbf{Local error statement.} Suppose an exact profile and a $C^3$ trial profile remain in a common region where $G$ is uniformly $L$-Lipschitz in $Y$. If their three initial data match and $|\widehat F'''+b/a|\le\varepsilon$ throughout an interval of length $\ell$, then Gronwall's inequality gives
\[
\|\widehat Y-Y\|_\infty\le\varepsilon(e^{L\ell}-1)/L,
\]
with $\varepsilon\ell$ at $L=0$. This classical estimate concerns the selected reduced IVP. A useful numerical value requires justified uniform bounds; sampled residuals alone do not provide them.

\newpage
\subsection{Semi-Analytical and Neural Approximation}
On a regular branch, calculate the Taylor polynomials $P_2$ and $P_4$ of $F$ at $\sigma_0$. The third derivative is $f=-b/a$ and the fourth is $f_\sigma+f_rv+f_vf$. For $z=\sigma-\sigma_0$, compare the two hard-constrained trial functions
\[
\widehat F_{\theta,N}(\sigma)=P_N(z)+z^{N+1}\mathcal N_\theta(z),
\qquad N=2,\ 4.
\]
Both satisfy the same initial position, slope and curvature; $N=4$ additionally embeds two ODE-derived derivatives. This changes both Taylor order and the remainder factor, so an improvement cannot be attributed to Taylor order alone. Taylor-plus-neural remainders are established prior art \cite{taylor}.

Train on the scaled reduced residual $R[\widehat F]/a$, using automatic derivatives of the single output $\widehat F$ \cite{pinn}. Monitor $R$ itself, $F'$ and $a$ separately so normalization cannot hide a failed physical or regularity check. A nonpositive $F'$ or a near-zero $a$ invalidates that evaluated branch. The network learns a profile, not a Lie algebra.

Compare standalone $P_2/P_4$, DOP853, and the two neural trials. Use identical two-hidden-layer, 16-unit tanh networks, three seeds, 128 training points and a distinct 1001-point evaluation grid. Fix the optimizer schedule: 5000 Adam steps followed by at most 200 L-BFGS iterations. Record actual function evaluations, CPU runtime and all failed runs; equal iterations need not mean equal cost.

Report profile and derivative errors, initial-data defects, reduced residuals, positivity and evaluation cost. On $\Omega=\{(t,s):0<t_{\min}\le t\le t_{\max},\ st^2\in I\}$, \eqref{eq:lift} gives $\|\mathcal E_3\|_{\infty,\Omega}\le t_{\min}^{-1}\|R\|_{\infty,I}$ for the profile interval $I=[1,1.5]$. This does not certify the Eulerian equations.

Symmetry-informed losses and transformed collocation points already exist \cite{liepinn,shah}. The proposed study instead fixes one verified invariant representation and tests its numerical reliability. Neural superiority, improved extrapolation and methodological novelty remain unproved.

\newpage
\subsection{Verification Pilot and Work Plan}
For $g=1$, $\beta=0.1$, use the source's cubic power-law family
\[
F_0(\sigma)=K\sigma^{1/3},\qquad
K^3=\frac{27g}{1-4g\beta}=45,\qquad 1\le\sigma\le1.5.
\]
Direct symbolic substitution verifies both this profile and the arbitrary-profile identity \eqref{eq:lift}. To avoid testing only a known closed form, also multiply $F_0''(1)$ by $1.01$, keeping $F_0(1)$ and $F_0'(1)$ fixed. This perturbs the profile data, \emph{not} the exponent $m$.

On 501 evaluation points, the classical solver's exact-control state error is $\LagrangeControlError$. For perturbed data, step/tolerance refinement differs by $\LagrangeRefinement$ and an independent Radau reference by $\LagrangeOracle$. Sampled minima are $F'=\LagrangeMinR$ and $|a|=\LagrangeMinA$. Refined finite differences check derivatives of the computed state curve independently of its ODE right-hand side.

For this perturbed profile, sampled $P_2$ and $P_4$ errors in $F$ are $\LagrangeTaylorTwo$ and $\LagrangeTaylorFour$, respectively. These are finite-interval diagnostics, not a Taylor convergence theorem or interval-certified error bounds. No PINN has been trained.

A separate exact traveling-wave control, motivated by Sections 3.6 and 3.9 of the source, satisfies the printed master equation but fails Eulerian Exner after reconstruction. It provides a negative model-equivalence test; details and runnable checks accompany the proposal.

\textbf{Work plan.} Weeks 1--3: source audit and derivation; 4--6: regularity bounds and classical checks; 7--10: neural experiments; 11--12: evaluation and writing. All computations remain on smooth compact regions away from $t=0$, $F'=0$ and $a=0$.

\section{Expected Outcomes}
Deliver an audited source-to-reduction derivation, a conditional reduced-profile error estimate, a reproducible benchmark and a measured Taylor--PINN comparison. The model-equivalence counterexample is retained as a limitation. The earlier Eulerian near-cubic pilot is historical evidence only, not validation of this proposal. Any originality claim requires further comparison with named prior work.

\newpage
\begin{thebibliography}{9}
\raggedright
\setlength{\itemsep}{3pt}
\setlength{\parskip}{0pt}
\bibitem{lagrange}
A. Paliathanasis, G. Leon, and P. G. L. Leach. Lie Symmetries and Similarity Solutions for a Shallow-Water Model with Bed Elevation in Lagrange Variables. \emph{Mathematics}, 14(3), 433, 2026.
\href{https://doi.org/10.3390/math14030433}{doi:10.3390/math14030433}.

\bibitem{olver}
P. J. Olver. \emph{Applications of Lie Groups to Differential Equations}. Second edition, Springer, 1993.
\href{https://doi.org/10.1007/978-1-4612-4350-2}{doi:10.1007/978-1-4612-4350-2}.

\bibitem{taylor}
C. K. Nagesh, S. Sankaranarayanan, R. Kaur, T. Sahai, and S. Jha. Taylor-Model Physics-Informed Neural Networks (PINNs) for Ordinary Differential Equations. \emph{PMLR}, 288, 621--642, 2025.
\href{https://proceedings.mlr.press/v288/nagesh25a.html}{proceedings.mlr.press/v288/nagesh25a.html}.

\bibitem{pinn}
M. Raissi, P. Perdikaris, and G. E. Karniadakis. Physics-informed neural networks: A deep learning framework for solving forward and inverse problems involving nonlinear partial differential equations. \emph{Journal of Computational Physics}, 378, 686--707, 2019.
\href{https://doi.org/10.1016/j.jcp.2018.10.045}{doi:10.1016/j.jcp.2018.10.045}.

\bibitem{liepinn}
T. Akhound-Sadegh, L. Perreault-Levasseur, J. Brandstetter, M. Welling, and S. Ravanbakhsh. Lie Point Symmetry and Physics-Informed Networks. \emph{NeurIPS}, 36, 42468--42481, 2023.
\href{https://arxiv.org/abs/2311.04293}{arXiv:2311.04293}.

\bibitem{shah}
A. H. Shah, N. R. Butt, A. Ahmad, and M. O. B. Saeed. Enhancing PINN Performance Through Lie Symmetry Group. Preprint, 2025.
\href{https://arxiv.org/abs/2509.26113}{arXiv:2509.26113}.
\end{thebibliography}
\end{document}
